% Image credits for the photographs used in this volume (Book 1).
% Machine-readable license records live in images/book1/CREDITS.md; keep
% the two files in sync when adding or removing a photograph.
\chapter*{Créditos das imagens}

Os desenhos deste livro são originais. As fotografias são utilizadas
sob as respectivas licenças livres, com agradecimento aos seus autores:

\begin{itemize}
  \item \emph{John Dalton} (capítulo sobre os átomos e as moléculas):
        gravura de Charles Turner a partir de Joseph Allen, 1834,
        Library of Congress via Wikimedia Commons, domínio público.
  \item \emph{A tabela dos elementos de Dalton, 1808} (capítulo sobre
        os átomos e as moléculas): prancha de \emph{A New System of
        Chemical Philosophy}, de John Dalton, digitalização do Internet
        Archive via Wikimedia Commons, domínio público.
  \item \emph{Bauxita} (capítulo sobre as matérias-primas): fotografia
        de Werner Schellmann via Wikimedia Commons, CC BY-SA 2.5.
  \item \emph{Granito} (capítulo sobre as substâncias puras):
        fotografia de Eurico Zimbres via Wikimedia Commons,
        CC BY-SA 2.0 br.
  \item \emph{Sulfato de cobre anidro} (capítulo sobre a identificação
        das substâncias): fotografia de W.~Oelen via Wikimedia Commons,
        CC BY-SA 3.0.
  \item \emph{Cristais de sulfato de cobre} (capítulo sobre a
        identificação das substâncias): fotografia de Stephanb via
        Wikimedia Commons, CC BY-SA 3.0.
  \item \emph{Antoine-Laurent Lavoisier e Marie-Anne Paulze Lavoisier}
        (capítulo sobre as equações balanceadas): pintura de
        Jacques-Louis David, 1788, Metropolitan Museum of Art via
        Wikimedia Commons, domínio público.
  \item \emph{Telefone de baquelite} (capítulo sobre os plásticos):
        fotografia de Holger Ellgaard via Wikimedia Commons,
        CC BY-SA 3.0.
  \item \emph{Ernest Rutherford} (capítulo sobre o interior do átomo):
        fotografia, George Grantham Bain Collection, Library of
        Congress, via Wikimedia Commons, domínio público.
  \item \emph{Halita} (capítulo sobre os íons): fotografia de Robert
        M.~Lavinsky via Wikimedia Commons, CC BY-SA 3.0.
  \item \emph{A Estátua da Liberdade} (capítulo sobre os metais e a
        corrosão): fotografia de Elcobbola via Wikimedia Commons,
        domínio público.
  \item \emph{Dmitri Mendeleiev} e \emph{a sua tabela de 1869}
        (capítulo sobre a tabela periódica): via Wikimedia Commons,
        domínio público.
  \item \emph{A Nebulosa do Caranguejo} (capítulo sobre os elementos):
        NASA, ESA, J.~Hester e A.~Loll, domínio público.
  \item \emph{Sódio} (capítulo sobre as camadas eletrônicas):
        fotografia de Dnn87 via Wikimedia Commons, CC BY-SA 3.0.
  \item \emph{Cloro numa ampola} (capítulo sobre as camadas
        eletrônicas): fotografia de W.~Oelen via Wikimedia Commons,
        CC BY-SA 3.0.
  \item \emph{Amedeo Avogadro} (capítulo sobre o mol): litografia a
        partir de um desenho de C.~Sentier, 1856, coleção Edgar Fahs
        Smith, via Wikimedia Commons, domínio público.
  \item \emph{Cristais de neve} (capítulo sobre a polaridade):
        fotografias de Wilson Bentley, 1902, via Wikimedia Commons,
        domínio público.
  \item \emph{Louis Pasteur} (capítulo sobre a estereoquímica):
        fotografia de Paul Nadar, via Wikimedia Commons, domínio
        público.
  \item \emph{Cristais de tartarato de amônio} (capítulo sobre a
        estereoquímica): fotografia de Marie-Lan Taÿ Pamart via
        Wikimedia Commons, CC BY 4.0.
  \item \emph{Svante Arrhenius} (capítulo sobre os ácidos e as bases):
        fotogravura, 1909, via Wikimedia Commons, domínio público.
  \item \emph{Uma pilha de Volta} (capítulo sobre as pilhas):
        fotografia de GuidoB via Wikimedia Commons, CC BY-SA 3.0.
  \item \emph{Wallace Carothers no seu laboratório} (capítulo sobre os
        polímeros): fotógrafo desconhecido, via Wikimedia Commons,
        domínio público.
\end{itemize}

Os links para as fontes completas e os endereços das licenças de cada
fotografia estão em \texttt{images/book1/CREDITS.md}, no repositório do
livro. Os pictogramas de perigo são os do Sistema Globalmente
Harmonizado de Classificação e Rotulagem de Produtos Químicos, tal como
desenhados para o pacote \texttt{ghsystem} do \TeX\ Live.

\bigskip
Ao lado das fotografias e dos desenhos originais, algumas figuras são
ilustrações geradas por inteligência artificial, produzidas com a
geração de imagens da OpenAI a partir de instruções escritas pelos
editores do livro e revisadas quanto à exatidão química antes da
inclusão. A instrução completa de cada imagem gerada está registrada em
\texttt{images/book1/ai/PROMPTS.md}, no repositório.
\clearpage
